\chapter{纯作业}
\section{第一周作业}
\subsection*{page44}

\begin{homework}[3]
设 $X_1, X_2, \cdots, X_n$ i.i.d., $X_1 \sim N(\mu,\sigma^2)$，其中 $\mu$ 是已知数，而 $\sigma^2$ 是未知的，指出下列各量哪些是统计量，哪些不是统计量，并说明理由：
\[
\sum_{i=1}^n X_i,\qquad
X_1+2\mu,\qquad
\max_{1\le i\le n}\{X_i\},\qquad
\sum_{i=1}^n \frac{(X_i-\mu)^2}{\sigma^2},\qquad
\sum_{i=1}^n (X_i-\mu)^2.
\]
\end{homework}

\begin{homework}[4]
某机床加工出的零件的直径服从正态分布 $N(\mu,\sigma^2)$，$\mu,\sigma^2$ 未知。为对 $\mu,\sigma^2$ 作出估计，随机选取加工出来的 $n$ 个零件进行测量，得到 $n$ 个数据
\[
x_1,x_2,\cdots,x_n.
\]
试写出该问题的统计模型（即参数空间和样本分布族）。
\end{homework}

\begin{homework}[6]
（柯西分布）设 $X_1,X_2,\cdots,X_n$ i.i.d.，$X_1$ 服从柯西分布，概率密度为
\[
f(x,\sigma)=\frac{1}{\pi\sigma}\frac{1}{1+(x/\sigma)^2},\qquad -\infty<x<\infty,
\]
求
\[
\overline{X}=\frac{1}{n}\sum_{i=1}^n X_i
\]
的抽样分布。
\end{homework}
\subsection*{page45}

\begin{homework}[1]
证明 $(2.2.3)$ 式。
\end{homework}

\begin{homework}[2]
（几何分布）设总体 $X$ 服从几何分布
\[
P\{X=k\}=p(1-p)^{k-1},\qquad k=1,2,\cdots,
\]
$X_1,X_2,\cdots,X_n$ 是抽自 $X$ 的简单随机样本，求 $X_{(1)},X_{(n)}$ 的分布。
\end{homework}

\begin{homework}[3]
（均匀分布）设 $X_1,X_2,\cdots,X_n$ i.i.d., $X_1\sim U\!\left[\theta-\frac12,\theta+\frac12\right]$，试分别求出 $X_{(1)}$ 和 $X_{(n)}$ 的概率密度。
\end{homework}
\section{第二周作业}
\subsection*{page45}

\begin{homework}[4]
设 $X_1,X_2,\cdots,X_n$ i.i.d., $X_1\sim F(x)$，$F(x)$ 有概率密度 $f(x)$，求：
\begin{itemize}
    \item[(1)] $X_{(1)}$ 和 $X_{(n)}$ 的联合分布函数及联合概率密度；
    \item[(2)] $R=X_{(n)}-X_{(1)}$ 的分布函数及概率密度。
\end{itemize}
\end{homework}

\begin{homework}[5]
从某总体抽取一个容量为 $5$ 的样本：
\[
-2.8,\ -1,\ 1.5,\ 2.1,\ 3.4.
\]
写出该样本的经验分布函数并画出它的图形。
\end{homework}

\begin{homework}[6]
证明：正态分布 $N(\mu,\sigma^2)$ 的上 $\alpha$（$0<\alpha<1$）分位数 $z_\alpha$ 与标准正态分布的上 $\alpha$ 分位数 $u_\alpha$ 之间有以下关系：
\[
z_\alpha=\mu+\sigma u_\alpha.
\]
\end{homework}
\subsection*{page46}

\begin{homework}[1]
（$\chi^2$ 分布）设 $X\sim \chi^2(n)$，证明：
\[
E(X^k)=2^k\left(\frac n2+k-1\right)\left(\frac n2+k-2\right)\cdots \frac n2,
\qquad k=1,2,\cdots.
\]
特别地，有
\[
E(X)=n,\qquad \operatorname{Var}(X)=2n.
\]
\end{homework}

\begin{homework}[2]
（$t$ 分布）设 $X\sim t(n)$，则当 $k<n$ 时，$E(X^k)$ 存在且有
\[
E(X^k)=
\begin{cases}
0, & \text{当 } k \text{ 为奇数；}\\[1em]
\dfrac{n^{k/2}\Gamma\!\left(\dfrac{k+1}{2}\right)\Gamma\!\left(\dfrac{n-k}{2}\right)}
{\sqrt{\pi}\,\Gamma\!\left(\dfrac n2\right)}, & \text{当 } k \text{ 为偶数。}
\end{cases}
\]
当 $k\ge n$ 时，$E(X^k)$ 是否存在？
\end{homework}

\begin{homework}[3]
（$\chi^2$ 分布）设 $X_1\sim \chi^2(m)$，$X_2\sim \chi^2(n)$，$X_1,X_2$ 相互独立，则 $X_1+X_2$ 与 $X_1/X_2$ 相互独立。
\end{homework}
\begin{homework}[5]
设 \(F\sim F(m,n)\)，证明
\[
F_{1-\alpha}(m,n)\,F_{\alpha}(n,m)=1,
\qquad \forall\,\alpha\in(0,1).
\]
\end{homework}
\begin{homework}
设 \(X_1,X_2,\cdots,X_n\) 独立，\(A_{2\times n}\) 已知，
\[
Y=AX,
\]
讨论 \(Y\) 与 \(X\) 的分布类型是否相同。
\end{homework}
\begin{homework}
（\(\chi^2\) 分布）设 \(X_1,X_2,\ldots,X_n\) i.i.d.，\(X_1\sim N(0,1)\)，
\[
\mathbf{X}=(X_1,\cdots,X_n)^{\mathrm T},
\]
\(A\) 为 \(n\) 阶对称方阵。证明：
\[
\mathbf{X}^{\mathrm T}A\mathbf{X}\sim \chi^2(r)
\]
当且仅当 \(A\) 为幂等矩阵，即
\[
A^2=A,
\]
并且
\[
\operatorname{rank}(A)=r.
\]
\end{homework}
\section{第三周}
\begin{homework}
设 \(X_1,\dots,X_n\) i.i.d., \(X_i\sim N(\mu,\sigma^2)\)，又设 \(X_{n+1}\sim N(\mu,\sigma^2)\)，且与 \(X_1,\dots,X_n\) 相互独立。记
\[
\bar X=\frac1n\sum_{i=1}^n X_i,\qquad
s^2=\frac1{n-1}\sum_{i=1}^n (X_i-\bar X)^2.
\]
证明
\[
T=\sqrt{\frac{n}{n+1}}\frac{\bar X-X_{n+1}}{s}\sim t(n-1).
\]
\end{homework}
\begin{homework}[2.4第1题]
设 \(X_1,\ldots,X_n\) i.i.d.，\(X_1\sim N(0,1)\)，
\[
s^2=\frac{1}{n-1}\sum_{i=1}^n (X_i-\overline{X})^2,
\]
求 \(\mathrm{Var}(s^2)\)。
\end{homework}

\begin{homework}[2.4第2题]
设 \(X_1,\ldots,X_m\) i.i.d.，\(X_1\sim N(\mu_1,\sigma^2)\)；\(Y_1,\ldots,Y_n\) i.i.d.，\(Y_1\sim N(\mu_2,\sigma^2)\)，且这两组随机变量相互独立。记
\[
s_1^2=\frac{1}{m-1}\sum_{i=1}^m (X_i-\overline{X})^2,\qquad
s_2^2=\frac{1}{n-1}\sum_{i=1}^n (Y_i-\overline{Y})^2.
\]
\begin{enumerate}
    \item 求
    \[
    \frac{(m-1)s_1^2+(n-1)s_2^2}{\sigma^2}
    \]
    的分布；
    \item 比较 \(s_1^2\) 与
    \[
    \frac{(m-1)s_1^2+(n-1)s_2^2}{m+n-2}
    \]
    的方差。
\end{enumerate}
\end{homework}

\begin{homework}[2.4第3题]
设 \(X_1,\ldots,X_n\) i.i.d.，\(X_1\sim N(0,1)\)，求
\[
Z=\frac{\sum_{i=1}^{n-1}X_i}{X_n}
\]
的分布。
\end{homework}
\begin{homework}[1]
（泊松分布）设 \(X_1,X_2,\cdots,X_n\) i.i.d., \(X_1\sim P(\lambda)\)，分别用一阶矩和二阶矩求 \(\lambda\) 的矩估计量。
\end{homework}

\begin{homework}[2]
（二项分布）设 \(X_1,X_2,\cdots,X_n\) i.i.d., \(X_1\sim B(k,p)\)，求 \(p\) 的矩估计量。
\end{homework}

\begin{homework}[3]
（均匀分布）设 \(X_1,X_2,\cdots,X_n\) i.i.d., \(X_1\sim U(\theta_1,\theta_2)\)，求 \(\theta_1,\theta_2\) 的矩估计量。
\end{homework}



\begin{homework}[5]
（双边指数分布）设 \(X_1,X_2,\cdots,X_n\) i.i.d., \(X_1\sim f(x)\)，
\[
f(x)=\frac{1}{2\sigma}\exp\!\left(-\frac{|x|}{\sigma}\right),
\]
求 \(\sigma\) 的矩估计量。
\end{homework}
\section{第五周作业}
\begin{homework}[7]
设 \(X_1,X_2,\cdots,X_n\) i.i.d.，试分别求下列总体中 \(\theta\) 的最大似然估计量。

\begin{enumerate}
    \setcounter{enumi}{1}
    \item （韦布尔分布）\(f(x;\theta)=\dfrac{x}{\theta^2}e^{-x^2/\theta^2},\ x>0,\ \theta>0.\)
\end{enumerate}
\end{homework}

\begin{homework}[8]
（均匀分布）设 \(X_1,X_2,\cdots,X_n\) i.i.d.，\(X_1\sim U\left(\theta-\dfrac{1}{2},\theta+\dfrac{1}{2}\right)\)，证明任意
\[
T\in \left[X_{(n)}-\frac{1}{2},\,X_{(1)}+\frac{1}{2}\right]
\]
都是 \(\theta\) 的最大似然估计量。
\end{homework}
\begin{homework}[10]
（双参数指数分布）设 \(X_1,X_2,\cdots,X_n\) i.i.d.，\(X_1\sim f(x;\alpha,\beta)\)，
\[
f(x;\alpha,\beta)=
\begin{cases}
\beta \exp\{-\beta(x-\alpha)\}, & x\ge \alpha,\\
0, & \text{其他},
\end{cases}
\]
其中 \(\beta>0\)。试求 \(\alpha,\beta\) 的最大似然估计量。
\end{homework}

\begin{homework}[11]
（截尾几何分布）设 \(X_1,X_2,\cdots,X_n\) i.i.d.，为来自截尾几何分布总体 \(X\) 的样本：
\[
P\{X=k\}=\theta^{k-1}(1-\theta),\qquad k=1,2,\cdots,r;
\]
\[
P\{X=r+1\}=\theta^r.
\]
记
\[
M=\#\{\,i: X_i=r+1,\ i=1,2,\cdots,n\,\},
\]
试证明 \(\theta\) 的 MLE 为
\[
\hat{\theta}
=
\frac{\sum_{i=1}^n X_i-n}{\sum_{i=1}^n X_i-M}.
\]
\end{homework}

\begin{homework}[12]
设 \(X_1,X_2,\cdots,X_n\) i.i.d.，\(X_1\sim f(x,\theta_1,\theta_2)\)，
\[
f(x,\theta_1,\theta_2)=
\begin{cases}
\dfrac{\theta_1}{\theta_2}x^{\theta_1-1}\exp\!\left(-\dfrac{x^{\theta_1}}{\theta_2}\right), & x\ge 0,\\[1ex]
0, & \text{其他},
\end{cases}
\]
其中 \(\theta_1\) 已知，求 \(\theta_2\) 的最大似然估计量。
\end{homework}
\begin{homework}[18]
设 \(X_1,X_2,\cdots,X_n\) i.i.d.，\(X_1\sim N(\mu,\sigma^2)\)，其中 \(\sigma^2\) 已知，\(\mu\) 的先验分布为 \(N(\nu,\tau^2)\)。

\begin{enumerate}
    \item 求 \(\mu\) 的后验分布；
    \item 求 \(\mu\) 的后验均值估计。
\end{enumerate}
\end{homework}

\begin{homework}[20]
（帕雷托（Pareto）分布）设 \(X_1,X_2,\cdots,X_n\) i.i.d.，\(X_1\sim U(0,\theta)\)，其中 \(\theta\) 的先验分布为 \(\mathrm{Pareto}(a,b)\)，其概率密度为
\[
\pi(\theta)=\frac{ba^b}{\theta^{b+1}},\qquad a<\theta<\infty,\ a>0,\ b>0.
\]

\begin{enumerate}
    \item 求 \(\theta\) 的后验分布；
    \item 求 \(\theta\) 的后验均值估计。
\end{enumerate}
\end{homework}
\begin{homework}[1]
（泊松分布）设 $X \sim P(\lambda)$，试求样本量为 $1$ 时 $g(\lambda) = e^{-2\lambda}$ 的无偏估计量，并说明该估计的不合理性。
\end{homework}

\begin{homework}[2]
（二项分布）设 $X \sim B(n,p)$，证明 $g(p) = \dfrac{1}{1+p^2}$ 不存在无偏估计量。
\end{homework}

\begin{homework}[3]
设 $X_1, X_2, \cdots, X_n$ i.i.d.，$X_1 \sim N(\mu, 1)$，证明 $g(\mu) = |\mu|$ 没有无偏估计量。
\end{homework}

\begin{homework}[4]
（均匀分布）设 $X_1, X_2, \cdots, X_n$ i.i.d.，$X_1 \sim U(\theta_1, \theta_2)$，求 $\theta_1, \theta_2$ 的无偏估计量。
\end{homework}

\begin{homework}[5]
设 $\hat{\theta}$ 和 $\theta^*$ 为 $\theta$ 的无偏估计量，若 $\tilde{\theta} = a\hat{\theta} + b\theta^*$，问：$a, b$ 满足什么条件时，$\tilde{\theta}$ 也是 $\theta$ 的无偏估计量？
\end{homework}

\begin{homework}[6]
设 $X \sim U(0,\theta)$，$X_1, X_2, \cdots, X_n$ i.i.d. 为来自 $X$ 的样本，试分别求 $\theta^2, \theta^{-1}$ 的无偏估计量。
\end{homework}
\section{第七周作业}
\begin{homework}[7]
设 $W$ 是 $\theta$ 的无偏估计量，若 $W$ 满足
\begin{enumerate}[(1)]
    \item $W$ 是观察值（样本）的线性函数；
    \item $E_\theta W = \theta, \forall \theta \in \Theta$；
    \item 若对任意一个满足 (1),(2) 的估计 $W^*$，有
    \[
    \mathrm{Var}_\theta(W) \leq \mathrm{Var}_\theta(W^*), \quad \forall \theta \in \Theta,
    \]
\end{enumerate}
则称 $W$ 是 $\theta$ 的最优线性无偏估计量。如总体 $X \sim N(\mu, \sigma^2)$，$(\mu, \sigma^2) \in \Theta = \{(\mu, \sigma^2) : -\infty < \mu < \infty, \sigma^2 > 0\}$。证明 $\overline{X} = \dfrac{1}{n}\sum_{i=1}^{n} X_i$ 是 $\mu$ 的最优线性无偏估计量。该结论只对正态总体正确吗？
\end{homework}

\begin{homework}[11]
（泊松分布）设 $X_1, X_2, \cdots, X_n$ i.i.d.，$X_1 \sim P(\lambda)$，证明：
\begin{enumerate}[(1)]
    \item $\overline{X}$ 是 $\lambda$ 的 UMVUE；
    \item $\tilde{\lambda}^2 = \dfrac{1}{n}\sum_{i=1}^{n} X_i^2 - \overline{X}$ 为 $\lambda^2$ 的无偏估计量。该估计是 UMVUE 吗？
\end{enumerate}
\end{homework}

\begin{homework}[12]
设 $X_1, X_2, \cdots, X_n$ i.i.d.，$X_1 \sim N(\mu, \sigma^2)$，已知 $\overline{X} = \dfrac{1}{n}\sum_{i=1}^{n} X_i$ 和 $s^2 = \dfrac{1}{n-1}\sum_{i=1}^{n}(X_i - \overline{X})^2$ 分别是 $\mu$ 和 $\sigma^2$ 的 UMVUE。计算 $\mu$ 和 $\sigma^2$ 的 C-R 下界。
\end{homework}

\begin{homework}[14]
在定理 3.3.3 的条件下，若 $\hat{g}(X_1, X_2, \cdots, X_n)$ 是 $g(\theta)$ 的估计量，满足
\[
E_\theta \hat{g}(X_1, X_2, \cdots, X_n) = g(\theta) + b(\theta), \quad \forall \theta \in \Theta \subset \mathbf{R},
\]
$b(\theta)$ 为估计的偏差，则有
\[
\mathrm{Var}_\theta(\hat{g}(X_1, X_2, \cdots, X_n)) \geq [b'(\theta) + g'(\theta)]^2 / nI(\theta).
\]
\end{homework}

\begin{homework}[15]
设 $\hat{g}_1(\boldsymbol{X})$ 和 $\hat{g}_2(\boldsymbol{X})$ 都是 $g(\theta)$ 的无偏估计量，并且 $\hat{g}_1(\boldsymbol{X})$ 还是 $g(\theta)$ 的 UMVUE，证明：$\hat{g}_1(\boldsymbol{X})$ 与 $\hat{g}_2(\boldsymbol{X})$ 的相关系数非负。
\end{homework}

\begin{homework}[16]
（0-1 分布）考虑例 3.2.12 中的贝叶斯估计量，其中 $\theta$ 的贝叶斯估计量为
$$
\hat{\theta} = \frac{\sum_{i=1}^n X_i + a}{n + a + b}.
$$
\begin{enumerate}[(1)]
    \item 求该估计量的均方误差。
\end{enumerate}
\end{homework}

\begin{homework}[1]
设 $X_n \xrightarrow{d} X$，$Y_n \xrightarrow{a.s.} a > 0$，证明：
\[
\frac{X_n}{Y_n} \xrightarrow{d} \frac{X}{a}.
\]
\end{homework}

\begin{homework}[1]
若估计序列 $\{\hat{\theta}_n\}$ 满足 $\lim_{n\to\infty} E_\theta(\hat{\theta}_n - \theta)^2 = 0$，$\forall \theta \in \Theta$，则称 $\hat{\theta}_n$ 为 $\theta$ 的均方相合估计量。证明：均方相合估计必是相合估计量。
\end{homework}

\begin{homework}[2]
（均匀分布）设 $X_1, X_2, \cdots, X_n$ i.i.d.，$X_1 \sim U(0, \theta)$，$\hat{\theta}_n = X_{(n)}$ 是 $\theta$ 的均方相合估计量吗？
\end{homework}

\begin{homework}[4]
（泊松分布）设 $X_1, X_2, \cdots, X_n$ i.i.d.，$X_1 \sim P(\lambda)$，证明：$\overline{X}$，$s_n^2$ 都是 $\lambda$ 的相合估计量。
\end{homework}

\begin{homework}[5]
（均匀分布）设 $X_1, X_2, \cdots, X_n$ i.i.d.，$X_1 \sim U(0,\theta)$，试证
\[
T(X_1, X_2, \cdots, X_n) = \sqrt[n]{X_1 X_2 \cdots X_n}
\]
为 $g(\theta) = \dfrac{\theta}{\mathrm{e}}$ 的相合估计量。
\end{homework}

\begin{homework}[6]
设总体的二阶矩存在，$X_1, X_2, \cdots, X_n$ 是总体的简单随机样本，则
\[
T(X_1, X_2, \cdots, X_n) = \frac{2}{n(n+1)} \sum_{i=1}^n i X_i
\]
是一阶原点矩的相合估计量。
\end{homework}

\begin{homework}[8]
（指数分布）设 $X_1, X_2, \cdots, X_n$ i.i.d.，$X_1 \sim f(x,\theta) = \dfrac{1}{\theta}\mathrm{e}^{-x/\theta}$，$x > 0$，试证 $\theta$ 的 MLE 之渐近分布为 $N\!\left(\theta, \dfrac{\theta^2}{n}\right)$。
\end{homework}